%%%%%%%%%%%%%%%%%%%%%%%%%%%%%%%%%%%%%%%%%%%%%%%%%%%%%%%%%%%%%%%%%%%%%%%%%%%%%%%
%% Seção 5 --- Conclusões
%%
%% Volta ao problema da Introdução, avalia as hipóteses H1 a H6 e não traz autores
%% nem citações novas, como pede o modelo do iCEV. Ao alterar a resposta ao problema,
%% alterar também o resumo e o abstract.
%%
%% MARCAÇÃO DE TRECHO EM ANDAMENTO --------------------------------------------
%% \emandamento{...} destaca em vermelho os trechos que mudam com as medições do
%% dispositivo montado. Ver nota no capitulo01.
%%%%%%%%%%%%%%%%%%%%%%%%%%%%%%%%%%%%%%%%%%%%%%%%%%%%%%%%%%%%%%%%%%%%%%%%%%%%%%%

\makeatletter
\@ifundefined{emandamento}{%
  \@ifundefined{textcolor}%
    {\newcommand{\emandamento}[1]{#1}}%
    {\newcommand{\emandamento}[1]{\textcolor{red}{#1}}}%
}{}
\makeatother

\chapter{Conclusões}
\label{cap:conclusoes}

Este trabalho perguntou quais modos de falha de motores uma cadeia de medição de baixo custo
permite detectar e quanto as decisões de estimação alteram o desempenho reportado. Em dado
real decimado para a taxa do dispositivo, com falso alarme medido de \SI{2,0 \pm 1,3}{\percent},
a detecção foi de \SI{75,0 \pm 3,5}{\percent} em desbalanceamento, de \SI{14,8 \pm 4,3}{\percent}
e \SI{20,7 \pm 7,0}{\percent} em defeito de rolamento puro, e de \SI{5,0 \pm 3,1}{\percent} e
\SI{7,2 \pm 4,3}{\percent} em desalinhamento horizontal e vertical. Com esse detector e esses
quatro descritores, a cadeia pode ser usada para detectar desbalanceamento, mas detecta pouco
defeito de rolamento sem desbalanceamento sobreposto.

As decisões de estimação mudaram os resultados. O limiar padrão marcou \SI{31,2 \pm
4,2}{\percent} da operação normal em dado real, contra \SI{2,0 \pm 1,3}{\percent} do limiar
calibrado. Para comparar dispositivos dessa classe, é preciso informar a banda, o limiar e o
protocolo de validação.

\section{Avaliação das hipóteses}
\label{sec:conc-hipoteses}

As hipóteses H1, H2, H4 e H5 não foram refutadas. A H3 não foi refutada em dado real, mas foi
refutada em um teste com sinal sintetizado. A H6 depende do dispositivo e ainda está em avaliação.

\textbf{H1, banda.} Não foi refutada. A detecção de desbalanceamento superou a de defeito de
        rolamento sem massa sobreposta, sob a mesma taxa de falso alarme. Esse resultado não
        isola o efeito da banda, porque o desalinhamento, dentro da banda, teve detecção ainda
        menor, e os descritores não medem a linha na frequência característica.
        \emandamento{Os experimentos mantêm conteúdo até \SI{500}{\hertz}, e o \emph{firmware} usa
        \SI{184}{\hertz}. A avaliação precisa ser refeita com essa banda.}

\textbf{H2, integração.} Não foi refutada. Contra uma senoide com integral conhecida, a soma
        acumulada errou $+\SI{73}{\percent}$ sem filtro passa-alta, $+\SI{156}{\percent}$ com ele
        e $+\SI{716}{\percent}$ sob viés, e a integração na frequência ficou a menos de
        \SI{0,01}{\percent} do valor analítico. Quase todo o erro da soma vem da constante de
        integração: com a média removida, ele cai para \SI{0,02}{\percent} sem filtro e para
        \SI{1,6}{\percent} com filtro.

\textbf{H3, limiar.} Em dado real, fora do treino, não foi refutada: o limiar padrão marcou
        \SI{31,2 \pm 4,2}{\percent} da operação normal, e o calibrado, \SI{2,0 \pm 1,3}{\percent}.
        Em sinal sintetizado, nas janelas do ajuste, os valores foram de \SI{42,1}{\percent} e
        \SI{1,4}{\percent}. Com um ensaio normal fora do treino, os dois marcaram
        \SI{100}{\percent}, o que refuta a hipótese nessa condição. Os escores foram idênticos nos
        cinco valores do parâmetro.

\textbf{H4, métrica.} Não foi refutada. Entre as cadeias de \SI{20}{\kilo\hertz} e de
        \SI{1}{\kilo\hertz}, a área global caiu \SI{3,7}{\percent}, e a área parcial caiu
        \SI{16,2}{\percent} com limite de \SI{10}{\percent} de falso alarme e \SI{20,3}{\percent}
        com limite de \SI{5}{\percent}. As duas cadeias usam conjuntos diferentes de descritores,
        e a queda não pode ser atribuída só à banda.

\textbf{H5, protocolo.} Não foi refutada. Um Random Forest supervisionado atingiu
        \SI{99,0 \pm 1,0}{\percent} de acurácia sob partição aleatória por janela e
        \SI{63,2 \pm 19,2}{\percent} sob partição por ensaio. O experimento não usou o detector do
        \emph{firmware}.

\textbf{H6, dispositivo.} \emandamento{Ainda em avaliação. Em bancada, o receptor aceitou \num{270}
        pacotes com CRC correto, e um pacote se perdeu. Em outra montagem, a taxa de amostragem ficou
        entre \SI{999,78}{\hertz} e \SI{999,88}{\hertz} em três ciclos, dentro do limite de
        \SI{2}{\percent}. Os critérios precisam ser repetidos na montagem final, com transmissor e
        receptor registrados no mesmo ensaio.}

\section{Limites dos resultados}
\label{sec:conc-alcance}

Os resultados de detecção vêm de dado real de uma bancada instrumentada, decimado em \emph{software}.
O sinal foi medido, mas a aquisição do dispositivo foi simulada. O treino usou apenas \num{49} arquivos
de operação normal, e todas as classes usam o mesmo canal da bancada. Essa bancada está abaixo do
limite de potência e de altura de eixo da norma, então os valores não podem ser estendidos diretamente
a máquinas industriais.

Os resultados de banda, integração, métrica e protocolo vêm de sinais sintetizados. Seus valores
absolutos dependem da amplitude do defeito, da frequência de ressonância e do nível de ruído escolhidos,
e devem ser lidos como ordens de grandeza. As relações entre os métodos devem depender menos desses
parâmetros, mas este trabalho não fez a varredura que confirmaria isso.

\emandamento{O dispositivo só foi ensaiado em bancada, sem motor e sem falha induzida, e a inferência
ainda não é executada nesse perfil. Portanto, os valores de detecção não são do dispositivo em operação.}

\section{Trabalho futuro}
\label{sec:conc-futuro}

\textbf{Projeto de \emph{hardware} próprio.} Este trabalho usou componentes em montagem de bancada. As
placas de circuito impresso, o invólucro, a bateria e o comportamento térmico ainda precisam ser
projetados e ensaiados. O invólucro e a fixação acrescentam uma ressonância entre o motor e o sensor, e
essa frequência deve ser medida por excitação de impacto.

\textbf{Ensaio em motor.} \emandamento{O dispositivo precisa ser ensaiado em um motor, com a inferência
ativada e com falhas induzidas. O piso de ruído também deve ser medido, para separar um defeito não
detectado de um defeito abaixo do ruído.}

\textbf{Varredura de taxa de amostragem.} A mesma gravação, decimada em várias taxas e analisada por um
único método e um único conjunto de descritores, permitiria separar o efeito da banda dos demais.

\textbf{Efeito do filtro em dado medido.} Em sinal sintetizado, o rebatimento mudou a curtose. Falta
medir o mesmo efeito em uma gravação real, com e sem filtro.

\textbf{Limiar em dado industrial.} A varredura do parâmetro de contaminação contra o falso alarme medido,
em vibração industrial de banda limitada, permitiria verificar se a calibração por risco alvo supera o
valor padrão.

\textbf{Treino no próprio microcontrolador.} A busca não encontrou taxa de falso alarme publicada para esse
caso, que é o de uma máquina sem histórico rotulado.

\textbf{Descritores de razão harmônica.} A detecção de desalinhamento foi a menor entre os modos medidos, e
os descritores adotados não medem a relação entre harmônicas. O teste pode usar as mesmas gravações
decimadas.
